\clearpage\thispagestyle{empty}
\pdfbookmark[0]{Quellen zu Juno (August 1811)}{juno1811aug-sources}
\begin{center}\textbf{Quellen zu Juno (August 1811)}\end{center}
\noindent Diese KI-Übertragung ins heutige Deutsch folgt dem Bericht in \href{https://gdz.sub.uni-goettingen.de/id/PPN236007467}{Gauß’ \textit{Werken}, Bd.~VI} (Göttingen, 1874), S.~330--331. Der Erstdruck erschien als Beitrag XXI, \textit{Fortgesetzte Nachrichten über den neuen Haupt-Planeten Juno}, in der \textit{Monatlichen Correspondenz}, Bd.~XXIV, August 1811, S.~\href{https://api.digitale-sammlungen.de/iiif/image/v2/bsb10538613_00196/full/1600,/0/default.jpg}{186}, \href{https://api.digitale-sammlungen.de/iiif/image/v2/bsb10538613_00197/full/1600,/0/default.jpg}{187} und \href{https://api.digitale-sammlungen.de/iiif/image/v2/bsb10538613_00198/full/1600,/0/default.jpg}{188}. Das digitalisierte Exemplar stammt aus der Bayerischen Staatsbibliothek.
\par\smallskip\noindent Der Erstdruck nennt für den Bericht vom September 1808 Bd.~XVIII, S.~270. Eine Fußnote auf S.~187 verweist für das Oppositionsergebnis auf Stück 92 der \textit{Göttingischen gelehrten Anzeigen}. Wachters Ephemeride für 1812 beginnt auf S.~188. Die Werkausgabe gibt die Daten in eckigen Klammern an. KI-Übertragung und KI-Satz.
\clearpage
